% SOURCE_AUTHORITY: sga5_fr_workpass.tex, lines 15235--15244 (LF-normalized unit).
% SOURCE_SHA256: 791F4EFFC5E02832D5D77ED03518C8156D6F07E4C8238B03545DB93D883FBB28
Para demostrar los lemas 1 y 2 necesitaremos el resultado siguiente:

\begin{statement}{Lema 3}
\begin{enumerate}[label=(\roman*)]
\item Sean $P$ y $Q$ dos $A$-módulos. Si $P$ es proyectivo, para toda aplicación $A_0$-lineal $f_0:T(P)\to T(Q)$ existe una aplicación $A$-lineal $f:P\to Q$ tal que $T(f)=f_0$.
\item Supongamos que el ideal $I$ es nilpotente. Para todo $A_0$-módulo proyectivo $P_0$ existe un $A$-módulo proyectivo $P$ tal que $T(P)\simeq P_0$. Si, además, $P'$ es un $A$-módulo plano tal que $T(P')\simeq P_0$, entonces se tiene $P'\simeq P$.
\end{enumerate}
\end{statement}

La afirmación (i) es clara, pues si $u_P$ (resp. $u_Q$) designa la aplicación canónica de $P$ sobre $T(P)$ (resp. de $Q$ sobre $T(Q)$), la compuesta $f_0\circ u_P$ es una aplicación $A$-lineal de $P$ en $T(Q)$ y $u_Q$ es sobreyectiva; como $P$ es proyectivo, existe una factorización de $f_0\circ u_P$ a través de $u_Q$: $f_0\circ u_P=u_Q\circ f$, donde $f$ es una aplicación $A$-lineal de $P$ en $Q$; entonces $T(f)=f_0$.
